\summarychapter{31}{Unwilling to Die for an Old Idea}{Denominations}

\begin{inbrief}
Critics call denominations a scandal, yet they remain the institutional hallmark of modern
Christianity. Christians are divided today partly because they are free to differ, as they were not
in earlier centuries. The chapter shows how the wars of religion wore out the old rule of one
religion for one territory, how the American colonies made room for many faiths, and how a group of
seventeenth-century Puritans worked out the denominational idea of the church.
\question{By what path did the denomination become the usual form of church life in the modern world?}
\end{inbrief}

Anyone reading the New Testament feels the distance between the church of the apostles, one temple of
God, and today's menagerie of sects, denominations and isms. It ought not to be so, yet denominations
will not soon disappear, because removing them would cost more than most people are willing to pay:
we are shocked by the fruit but will not lay the ax to the root. One clear difference between the
sixteenth century and the seventeenth is the acceptance of religious differences. ``I disapprove of
what you say, but I will defend to the death your right to say it'' is a sentiment most Christians
now accept, and they accept it as modern people rather than as Christians.

\ssection{Suppression of the Nonconformist}
The saying is probably Voltaire's (1694--1778). Neither Luther nor Loyola could have said it. In the
Reformation dissent was neither a virtue nor a right, and Reformers were as eager as Catholics to
suppress it, for both believed that Christian truth holds society together and that only one side
has it. That God's Word might lie on either side of a battle line gained a hearing only when both
sides were exhausted. In Germany Lutheran princes fought the Catholic emperor to a stalemate, and the
Peace of Augsburg (1555) adopted the territorial principle: the ruler of each region decides whether
his subjects are Lutheran or Catholic. This offshoot of the idea of the Christian prince guaranteed
some people their worship and also the prince's right to persecute those who disagreed. It foretold
the death of Christendom, for if one prince can impose one religion and another, a dozen miles away,
a different one, the norm of Christian truth is shattered.

Events ran ahead of ideas, and more wars followed on the old principle. From 1562 to 1598 France
suffered civil wars between Catholics and Huguenots, ended in fatigue by the Edict of Nantes (1598),
which gave the Huguenots religious freedom and political control of parts of the country while
Catholicism stayed the official religion. Between 1560 and 1618 the Calvinist Dutch won independence
from Catholic Spain, while the southern provinces, today's Belgium, remained Catholic and Spanish
much longer.

\ssection{The Thirty Years' War}
These were the prelude to the last and worst of the wars of religion (1618--1648). At the start religion was the cause and politics the accompaniment; at the end it was a savage contest for power in which religion was only the accompaniment. That makes it an apt sign of the passage from the age of the Reformation to that of reason and revival. Augsburg had ignored the Calvinists, and early in the seventeenth
century Protestant and Catholic princes formed rival leagues. Fighting broke out in 1618. Ferdinand
II, Jesuit-educated, king of Bohemia and soon Holy Roman Emperor, tried to uproot Protestantism in
Bohemia. The Protestant nobles revolted and offered the crown to the Calvinist Frederick V of the
Palatinate. In 1620, near Prague, the imperial forces, charging to the cry ``For the Virgin Mary,''
crushed the Bohemians, confiscated the rebels' estates and gave the Jesuits the University of Prague,
once the home of Hus.

Christian IV of Lutheran Denmark then entered the war, eager also for German territory, and was
routed in the Harz Mountains in 1626. Next came Gustavus Adolphus of Sweden, moved by deep conviction
and by the emperor's hold on the Baltic shore. The ``Lion of the North'' won a series of victories
that carried him as far as Munich, but at L\"utzen in 1632, though his army won, he was killed. The
war dragged on with the outcome already plain: Catholics could not subdue the Protestant north nor
Protestants the Catholic south. In its last years religion hardly mattered, as France and Spain, both
nominally Catholic, fought over the Rhineland.

Germany lay in ruins, divided into three hundred independent states, and the emperor's dream of
revived authority was gone. Out of sheer futility zeal cooled and people began to question the
territorial idea. The Peace of Westphalia (1648) recognized Calvinism alongside Lutheranism and
Catholicism, allowed princes for the first time to tolerate both Protestants and Catholics in their
lands, and shut the pope out of German religious affairs. Innocent X condemned the treaty and was
ignored by both sides; after more than a thousand years the state did its affairs as if there were no pope. The real victor was the nation-state. Christians in the West would now live under
governments whose citizens belonged to various churches. Some states came to require a general
belief in God for office or university, while being Catholic or Presbyterian became a private matter.

\ssection{The New Way in America}
Across the Atlantic the English colonists set out to build distinctly religious societies. In time
the aim of a holy city that nurtures holy people became a city where people could be holy, for new
arrivals, an indifferent second generation and the open frontier all worked against uniformity.
Permanent settlement began at Jamestown, Virginia, in 1607. The Pilgrims landed in 1620, and about twenty-five thousand Puritans followed to New England from 1629 to 1642. Outside Virginia the
English authorities did not try to impose religious uniformity. The colonies were commercial
ventures that needed settlers, and the promise of toleration was a strong inducement. Quakers came to
Pennsylvania, Catholics to Maryland, Dutch Reformed to New York, and later Swedish Lutherans, French
Huguenots, English Baptists and Scottish Presbyterians.

The great exception was Congregational Massachusetts Bay, which intended a new Zion with ``a due
form of government both civil and ecclesiastical,'' in Governor John Winthrop's words. Its charter
happened to omit the clause requiring the company's headquarters to be in England, so the saints
took the charter with them and ran what was in effect an independent republic. For two generations
they enforced conformity, punishing absence from church, denial of Christ's resurrection or
irreverence toward the Bible. Yet Puritanism carried its own impulses toward toleration. One was its
devotion to Scripture. Pastor John Robinson had told the Pilgrims that ``the Lord hath more truth yet
to break forth out of his holy Word,'' and New England kept producing dissenters who appealed to
what the Bible had shown them, such as Roger Williams and Anne Hutchinson. Another was the
wilderness. Dissenters did not need to go underground; they could move on, across the river or over
the mountain. With English policy these undercut Puritan intolerance. The colony lost its charter in
1684, and in 1687 the royal governor seized Boston's Old South meeting house for Anglican worship.

\ssection{The Idea of Denominations}
This diversity required a new understanding of the church. The word \emph{denomination} came into use
about 1740, in the early years of the revival led by Wesley and Whitefield, but the theory was worked
out a century earlier. In its original sense it is the reverse of sectarianism. A sect holds that Christ's authority belongs to it and to no one else, and so shuts others out by definition. A denomination is an inclusive
term: the group called, or denominated, by a particular name is one member of a larger body, the church, in which every denomination has a place. No one structure can claim to be the true church; each denomination is one form, in worship and organization, of the church's wider life.

The Reformers planted the seed. Luther held that the true succession is of believers, not bishops,
and that the outward forms of the church should give the Word free course rather than block it;
Calvin wrote that no one can draw the precise boundaries of the church or number the elect. But
neither followed this through, and both suppressed dissent. The real architects were the
Independents, or Congregationalists, the minority at the Westminster Assembly, where the Presbyterian
majority wrote the Confession and Catechisms. Anxious not to divide ``the godly Protestant party,''
these Dissenting Brethren stated four truths. Because human beings do not always see the truth
clearly, disagreement over the church's outward form cannot be avoided. Such differences, though
not about fundamentals, are not matters of indifference, since every Christian must practice what he
believes the Bible teaches. Since no church understands God's truth completely and finally, a single structure can never stand for the whole of the true church. And separation is not in itself schism:
Christians may differ on many matters and remain one in Christ.

The theory thus sought unity in an inward religious experience and allowed diversity in its outward
expression. It did not spring from indifference to doctrine. The Independents limited the one true church to people who agreed about the essentials of the faith. In England the idea
found limited acceptance, for the Church of England kept its favored place even after the Act of
Toleration (1689) gave Presbyterians, Congregationalists, Baptists and Quakers freedom of worship. In
the American colonies it won growing acceptance as God's answer to multiplying faiths. Its authors
never imagined the hundreds of groups now sheltered under the word and did not mean to reduce
Christianity to vague sincerity. They only knew that bigotry and bloodshed in Christ's name were not
the way forward. The denominational form has marked recent centuries, Shelley concludes, since nothing better has yet been proposed, and not because it is the ideal.

\begin{keyfacts}
\keypeople{Voltaire; Ferdinand II; Frederick V of the Palatinate; Christian IV of Denmark; Gustavus
Adolphus; Innocent X; John Winthrop; John Robinson; Roger Williams; Anne Hutchinson; the Dissenting
Brethren of Westminster.}
\keydates{%
  \kd{1555}{Peace of Augsburg: each ruler decides the religion of his territory}
  \kd{1598}{Edict of Nantes grants the Huguenots freedom}
  \kd{1607}{Jamestown, Virginia, the first permanent English settlement}
  \kd{1618--1648}{Thirty Years' War}
  \kd{1632}{Gustavus Adolphus killed at L\"utzen}
  \kd{1648}{Peace of Westphalia: Calvinism recognized; the pope excluded from German affairs}
  \kd{1684}{Massachusetts Bay loses its charter}
  \kd{1689}{Act of Toleration in England}
  \kd{c. 1740}{The word ``denomination'' comes into use}}
\keyterms{Denomination and sect; territorial principle; wars of religion; nation-state; toleration;
Independents (Congregationalists); Dissenting Brethren; denominational theory of the church.}
\end{keyfacts}
